\chapter*{主要缩略语表}
\ifdefined\HCode\else\addcontentsline{toc}{chapter}{主要缩略语表}\fi
\pagestyle{mystyle}

\small

% 分类小标题行：PDF 用 \multicolumn 跨三列；HTML(tex4ht) 下 \multicolumn
% 作为整行会触发 “Misplaced \omit”，故改为普通三列行（首列放分类名）。
\ifdefined\HCode
  \newcommand{\catrow}[1]{\textbf{\textit{#1}} & & \\}
\else
  \newcommand{\catrow}[1]{\multicolumn{3}{l}{\textit{#1}} \\}
\fi

% tex4ht(HTML) 与 PDF 用不同表格环境：
%   PDF  → longtable（跨页、重复表头）+ booktabs 横线
%   HTML → 普通 tabular + \hline（tex4ht 对 longtable 的 \endhead 及
%          booktabs 的 \noalign 横线支持差，会报 Misplaced \omit/\cr
%          并污染 xetex .xdv 导致 postamble 致命错误）
\ifdefined\HCode
  % HTML 分支内：把 booktabs 横线降级为 \hline，\addlinespace 置空
  \renewcommand{\addlinespace}[1][]{}%
  \begin{center}
  \begin{tabular}{p{2.5cm}p{6.8cm}p{4.2cm}}
  \hline
  \textbf{缩略语} & \textbf{英文全称} & \textbf{中文含义} \\
  \hline
\else
\begin{longtable}{p{2.5cm}p{6.8cm}p{4.2cm}}

\toprule
\textbf{缩略语} & \textbf{英文全称} & \textbf{中文含义} \\
\midrule
\endfirsthead

\toprule
\textbf{缩略语} & \textbf{英文全称} & \textbf{中文含义} \\
\midrule
\endhead

\midrule
\multicolumn{3}{r}{\small\itshape 续下页} \\
\endfoot

\bottomrule
\endlastfoot
\fi

%================================================
\catrow{引力波探测器与观测任务}
\addlinespace[0.3em]

LIGO & Laser Interferometer Gravitational-Wave Observatory & 激光干涉引力波天文台 \\
LISA & Laser Interferometer Space Antenna & 激光干涉仪空间天线 \\
KAGRA & Kamioka Gravitational Wave Detector & 神冈引力波探测器 \\
ET & Einstein Telescope & 爱因斯坦望远镜 \\
CE & Cosmic Explorer & 宇宙探索者 \\
PTA & Pulsar Timing Array & 脉冲星计时阵列 \\
TDI & Time Delay Interferometry & 时间延迟干涉测量 \\

\addlinespace[0.6em]

%================================================
\catrow{引力波源与天体系统}
\addlinespace[0.3em]

BBH & Binary Black Hole & 双黑洞系统 \\
BNS & Binary Neutron Star & 双中子星系统 \\
NSBH & Neutron Star--Black Hole & 中子星--黑洞系统 \\
MBHB & Massive Black Hole Binary & 大质量双黑洞系统 \\
EMRI & Extreme Mass Ratio Inspiral & 极端质量比旋入 \\
SGWB & Stochastic Gravitational-Wave Background & 随机引力波背景 \\

\addlinespace[0.6em]

%================================================
\catrow{信号处理、统计与数值方法}
\addlinespace[0.3em]

PSD & Power Spectral Density & 功率谱密度 \\
SNR & Signal-to-Noise Ratio & 信噪比 \\
FAR & False Alarm Rate & 误报率 \\
ROC & Receiver Operating Characteristic & 接收者工作特征曲线 \\
AUC & Area Under the Curve & 曲线下面积 \\
FF & Fitting Factor & 匹配度 \\
DQ & Data Quality & 数据质量 \\
MCMC & Markov Chain Monte Carlo & 马尔可夫链蒙特卡洛方法 \\
FFT & Fast Fourier Transform & 快速傅里叶变换 \\
SVD & Singular Value Decomposition & 奇异值分解 \\

\addlinespace[0.6em]

%================================================
\catrow{理论物理与波形建模}
\addlinespace[0.3em]

GR & General Relativity & 广义相对论 \\
PN & Post-Newtonian & 后牛顿近似 \\
NR & Numerical Relativity & 数值相对论 \\
SPA & Stationary Phase Approximation & 驻相近似 \\
TT & Transverse-Traceless & 横向无迹规范 \\

\addlinespace[0.6em]

%================================================
\catrow{机器学习与人工智能方法}
\addlinespace[0.3em]

CNN & Convolutional Neural Network & 卷积神经网络 \\
MF-CNN & Matched-Filtering Convolutional Neural Network & 匹配滤波卷积神经网络 \\
RNN & Recurrent Neural Network & 循环神经网络 \\
LSTM & Long Short-Term Memory & 长短期记忆网络 \\
CVAE & Conditional Variational Autoencoder & 条件变分自编码器 \\
CNF & Continuous Normalizing Flow & 连续归一化流 \\
SBI & Simulation-Based Inference & 基于模拟的推断 \\
NPE & Neural Posterior Estimation & 神经后验估计 \\
NLE & Neural Likelihood Estimation & 神经似然估计 \\
NRE & Neural Ratio Estimation & 神经比率估计 \\
LLM & Large Language Model & 大语言模型 \\
MCTS & Monte Carlo Tree Search & 蒙特卡洛树搜索 \\
AHD & Automated Heuristic Design & 自动启发式设计 \\
PINN & Physics-Informed Neural Network & 物理信息神经网络 \\

\addlinespace[0.6em]

%================================================
\catrow{数据集、挑战赛与观测运行}
\addlinespace[0.3em]

GWTC & Gravitational-Wave Transient Catalog & 引力波暂现源目录 \\
LDC & LISA Data Challenge & LISA 数据挑战赛 \\
TDC & Taiji Data Challenge & 太极数据挑战赛 \\
MLGWSC-1 & 1st ML Gravitational-Wave Search Challenge & 第一届机器学习引力波搜索挑战赛 \\
O1/O2/O3/O4 & Observing Run 1/2/3/4 & 第1/2/3/4轮科学观测 \\

\ifdefined\HCode
  \hline
  \end{tabular}
  \end{center}
\else
\end{longtable}
\fi
